% Solutions to the Chapter 11 exercises that are not code, from lectures/L11/appendix/b_exercises.md.
% Exercises 11.6 and 11.7 are Code and are not answered; see chapter.tex for why.

\section{Chapter 11: Kernel Frameworks}
\label{sol:sec:c11}
Answers to the exercises in \secref{c11:sec:exercises}. Exercises~\ref{c11:ex:6} and~\ref{c11:ex:7}
are Code, and are checked on the target: the first as its Check yourself line says, the second with
the same tools that \file{run.sh} uses.

% ------------------------------------------------------------------------------------------
\solution[fillaccent2]{c11:ex:1}{What a framework gives you}{Recall}

\xpart{a} The five of \secref{c11:app:a1}: no way for a program that does not already know about it
to find it; no documented interface, only the one in your head and in the source; no existing
tools, so every user writes their own reader; no power management, so suspend and resume do
nothing; and no way to say what the device \emph{is}, only what it does. IIO would also have given
it a scale and units, and a buffered stream to userspace (\secref{c11:app:a8}).

\xpart{b} You implement an interface somebody else designed, and in exchange you stop having to
design, document, version and support one of your own, and every tool that already speaks it works
with your device.

\xpart{c} First, it is an ABI as soon as a second program uses it: its numbers and structure
layouts are frozen, and you may add to them but never change or remove anything. Second, it is a
portability trap: a 32-bit process on a 64-bit kernel lays the argument structure out differently,
so the \code{ioctl} needs a \code{compat_ioctl} handler, or a structure with fixed-width types and no
implicit padding. Third, it is undiscoverable: nothing lists the devices that implement it, and no
tool speaks it. The second is the one that appears only on a particular kind of system: a 64-bit
kernel running 32-bit userspace, which is common in embedded work, an arm64 kernel under an older
32-bit Arm root filesystem for instance. It is invisible until somebody runs a 32-bit binary.

\xpart{d} There is none, which is the point. Nothing records which drivers implement which
\code{ioctl} numbers: a number means something only to the driver that defined it, and the only way
to find the devices is to know their names already. Trying a device with an \code{ioctl} it does
not implement fails with \code{ENOTTY} at best, and does something else entirely if its driver
happened to pick the same number.

% ------------------------------------------------------------------------------------------
\solution[fillaccent2]{c11:ex:2}{The tour}{Recall}

\xpart{a} \textbf{i2c} and \textbf{iio}: an I2C client, because that is how it is attached, and an
IIO device, because it measures. A temperature sensor that watches the health of the board itself
is often a \code{hwmon} device instead, a subsystem the tour in \secref{c11:app:a3} leaves out.

\xpart{b} \textbf{gpio}. On a real board the LED is usually handed on to the \code{leds} subsystem
through the \code{gpio-leds} driver, which consumes the line; the pin is still a GPIO.

\xpart{c} \textbf{watchdog}.

\xpart{d} \textbf{input}: it is a source of human input. Its driver usually reads the encoder's two
lines through gpio, but what the device is, is input.

\xpart{e} \textbf{mtd}: raw flash, with no controller in front of it doing wear levelling or
bad-block handling.

\xpart{f} \textbf{spi} and \textbf{iio}.

\xpart{g} \textbf{clk}: a clock source, described once in the device tree and shared by its three
consumers, whose enables the clock framework counts.

\xpart{h} a) and f), each an \textbf{i2c} or \textbf{spi} device and an \textbf{iio} device. That
is not a contradiction, because the two answer different questions, the same two that
\file{/sys/bus} and \file{/sys/class} answer (\secref{c10:app:a2}): the bus is how the device is
attached and how a driver is matched to it, and the class is what the device is to userspace. One
driver does both. It registers with the bus as an I2C or SPI driver, and its \code{probe} registers
an IIO device.

\xpart{i} \textbf{regmap}. It is a driver-side library rather than an interface: it gives a driver
register access over I2C, SPI, MMIO and other buses from a description, the register width and
which registers are readable, writable, volatile or cacheable, and does the transfers, the byte
order, the read-modify-write and the caching itself. The only view of it from outside is in
debugfs, which is not an interface.

% ------------------------------------------------------------------------------------------
\solution[fillaccent3]{c11:ex:3}{Reading an IIO device}{Hand calculation}

\xpart{a} With no offset the value is the raw count times the scale:
$2048 \times 0.805664 = 1649.999872$, so \textbf{1650.0\,mV}, or 1.650\,V. The units of the scale
are millivolts per count. The IIO ABI says so: the documentation of \code{in_voltageY_raw} states
that its units after the scale and offset are applied are millivolts
(\file{Documentation/ABI/testing/sysfs-bus-iio}:141--153), so the scale is whatever turns one count
into millivolts. The number gives the device away, too: $3300 / 4096 = 0.8056640625$, a 12-bit
converter with a 3.3\,V reference, and 2048 is half of full scale.

\xpart{b} The kernel does not use floating point, so a driver cannot cheaply produce a fractional
voltage, while the raw count is exactly what the hardware delivers and the scale can be stated
exactly, as a fraction if need be. In buffered mode the samples go to userspace as the hardware
produced them, at full rate, with no conversion per sample in the kernel, and userspace applies
one scale to all of them. And the scale belongs to the channel's configuration, which can change
with its gain or range, so it is published once rather than folded into every value.

\xpart{c} The offset is added to the raw value before scaling
(\file{Documentation/ABI/testing/sysfs-bus-iio}:455--457):
$(2048 + (-100)) \times 0.805664 = 1948 \times 0.805664 = 1569.433472$, so \textbf{1569.4\,mV},
about 1.569\,V. The order is raw, plus offset, times scale.

\xpart{d} Only \code{raw}. A counter has no unit to convert to, so a \code{scale} would claim one and
an \code{offset} would claim a zero point; without them userspace takes the value as a plain count.
\code{type} should say what the value is, and \code{IIO_VOLTAGE} tells every tool that reads it that
it is a voltage, which is false. IIO does have \code{IIO_COUNT}, which makes \code{in_count0_raw}
(\file{include/uapi/linux/iio/types.h}:44), but the ABI marks that attribute as deprecated in favour
of the counter subsystem (\file{Documentation/ABI/testing/sysfs-bus-iio}:2005--2009), so a device
that really is a counter belongs there, not in IIO. The lab uses \code{IIO_VOLTAGE} because
\code{qa-dev} stands in for a converter, and its counter mode is a test pattern rather than what the
device is.

% ------------------------------------------------------------------------------------------
\solution[fillaccent4]{c11:ex:4}{The pin that is not yet a GPIO}{Design}

\xpart{a} \textbf{pinctrl}, and the configuration is written in the device tree. The pin
controller's node holds groups naming the pins and the function each is to have, with their pull
and drive settings, and the node of the device that uses them selects one by phandle:

\begin{console}
pinctrl-names = "default";
pinctrl-0 = <&uart0_pins>;
\end{console}

\xpart{b} No. The driver core has the pin controller apply the device's \code{default} state before
\code{probe} is called (\code{pinctrl_bind_pins} in \code{really_probe},
\file{drivers/base/dd.c}:635), so by the time your code runs the pins are already set. Which
function a pin has depends on the board, the same pin of the same SoC being a UART line on one
board and a GPIO on the next, so it belongs in that board's tree. And the multiplexing registers
belong to the pin controller, shared by every device on the SoC: a driver that writes them
hardcodes one SoC's register layout and one board's wiring, and races every other user of the same
register.

\xpart{c} \code{devm_clk_get(&pdev->dev, NULL)} to get the clock, checked with \code{IS_ERR}, and
\code{clk_prepare_enable(clk)} to turn it on, checked for a negative return;
\file{include/linux/clk.h} also has \code{devm_clk_get_enabled}, which does both and turns the clock
off again at unbind. Writing the gate register directly instead goes wrong because the clock
framework counts enables, turns a clock off when the count reaches zero, and does not know about
you. When the other driver disables the clock, in its suspend or its \code{remove}, the count
reaches zero and the framework gates it, and your device stops responding underneath you. The
other way round, if your \code{remove} writes the gate off, the other driver's device dies. Even
with no second driver, the framework switches off every clock nobody has claimed at the end of boot
(\code{clk_disable_unused}, \file{drivers/clk/clk.c}:1531--1570), yours included; and it enables a
clock's parents before the clock, which a write to one gate does not.

\xpart{d} Get it with \code{devm_regulator_get} and turn it on with \code{regulator_enable}. The
framework counts its consumers and switches the supply off only when the last one lets go; it
enables the supply's own parent supplies first; it holds every consumer to the voltage limits and
permissions the board declares in its device tree; it waits for the supply to settle before
returning; and once boot has finished it switches off every regulator no consumer has claimed
(\code{regulator_late_cleanup}, \file{drivers/regulator/core.c}:6336--6361), which would include
yours if you had only written the register. The write itself would also mean talking to the power
management chip over its bus, which belongs to that chip's driver.

\xpart{e} \code{-EPROBE_DEFER}. The lookups return it themselves when the provider has not
registered yet: \code{devm_clk_get}, \code{devm_regulator_get}, and the pin controller lookup,
which the core does before \code{probe} even runs. \code{probe} passes it on, ideally through
\code{dev_err_probe}, which records why. The core puts the device on its list of deferred devices
and tries it again whenever another driver binds successfully (\file{drivers/base/dd.c}:36--50), so
it probes once its providers have; anything still waiting at the end of boot is listed in
\file{/sys/kernel/debug/devices_deferred}.

% ------------------------------------------------------------------------------------------
\solution[fillaccent4]{c11:ex:5}{When a character device is right}{Design}

\xpart{a} A block in an FPGA on a company's own board that runs a protocol of the product's own:
say, a controller that takes a list of timed commands for a production machine's actuators and
reports as each one completes. It fits no subsystem because it is not a member of any class of
hardware that other vendors make. It does not measure a value, it is not a set of pins, a network
interface or a store, and its operations are defined by the product, so no subsystem's model has a
place for them. An answer that names a subsystem instead has to argue that the device's operations
are that subsystem's: if the block turned out to be a set of PWM outputs, it would belong in
\code{pwm}.

\xpart{b} \code{misc_register}. One device needs one node, and the misc device gives it one: a
dynamic minor under major 10 (\file{include/linux/miscdevice.h}:74,
\file{include/uapi/linux/major.h}:26), with the class device and so the node created for it, in one
call undone by one call. Chapter~\ref{c5:ch}'s long way, a number, a \code{cdev}, a class and a
device, each with its unwind, is for a driver that needs a range of minors or a major of its own,
such as many nodes for one device.

\xpart{c} A specification of each: its number, built with \code{_IOR}, \code{_IOW} or \code{_IOWR}
with the right direction and size; the argument's layout and what every field means; which errors
it returns and when; and the header, in the kernel's userspace API form. A layout that is the same
for a 32-bit and a 64-bit caller, with fixed-width types, explicit padding and any pointer carried
as a \code{__u64}, or else a \code{compat_ioctl} that translates. Room to grow, a flags field that
must be zero today or a size field, so that the next change does not need a new number. And the
promise never to change or remove any of it. For how long: for as long as any program that uses it
exists, which in practice means forever, because there is no deprecation period for a customer's
binary (\secref{c11:app:a2}).

\xpart{d} \textbf{The \code{long}} is 4 bytes in a 32-bit process and 8 in a 64-bit kernel, so the
structure has a different size and layout for the two; and because \code{_IOR} encodes the size in
the command number, a 32-bit caller's number does not even match the 64-bit kernel's. It needs a
\code{compat_ioctl} that knows both layouts. \textbf{The pointer} has the same size problem, and a
worse one: a pointer means nothing on the other side of the boundary. A kernel pointer handed to
userspace leaks a kernel address, which defeats address-space randomisation, and is of no use to
the program; a userspace pointer has to be carried as a \code{__u64} and turned back with
\code{u64_to_user_ptr} (\file{include/linux/kernel.h}:50). The fix for both is fixed-width types.

\xpart{e} ``No subsystem fits'' is a statement about the device, and it is a reason: there is no
existing interface, so you write one and accept that you own it. ``Implementing a subsystem
interface looked like more work'' is a statement about the author's afternoon, and it is an excuse:
the interface is a cost paid once, by you, and the invented ABI a cost paid for as long as the
device is in use, by you and by everyone who has to write a tool for it.

% ------------------------------------------------------------------------------------------
\solution[fillaccent4]{c11:ex:8}{What you gave up}{Design}

\xpart{a} No. A read of \code{in_voltage0_raw} returns at once with the most recent sample, whether
or not a new one has arrived since the last read; nothing makes it wait, and \code{poll} on the
attribute is not woken by new samples, because the driver never tells sysfs about them. What has
been lost is waiting for data without spinning, and with it every sample but the latest: whatever
arrives between two reads is gone, and so are the order and the timing of the samples.

\xpart{b} IIO's buffers, described in \file{Documentation/driver-api/iio/buffers.rst} and
\file{triggered-buffers.rst}: the driver pushes each sample into a buffer in the kernel as it
arrives, and userspace enables the channels it wants and reads whole batches of samples from
\file{/dev/iio:deviceN}, blocking or polling until they are there. A driver must add a
\code{scan_index} and a \code{scan_type} to its channel and set up a buffer, which for a device
that raises its own interrupt per sample, as \code{qa-dev} does, is \code{devm_iio_kfifo_buffer_setup}
with the handler calling \code{iio_push_to_buffers} (a device that has to be told when to sample
uses \code{devm_iio_triggered_buffer_setup} and a trigger instead), and it needs
\code{CONFIG_IIO_BUFFER} and \code{CONFIG_IIO_KFIFO_BUF}, which \file{kernel/qa.config} already
enables.

\xpart{c} Anything that is not a channel value. The interrupt timestamp in \code{TS_LO} and
\code{TS_HI}, which Chapter~\ref{c12:ch} is built on, or the overrun flag in \code{STATUS}: with a
character device of your own you added a field to what \code{read} returned, or an \code{ioctl}. In
IIO you have to find the framework's place for it, a timestamp channel or an event, or add an
attribute of your own, which is new ABI again. A smaller one on this target: the module now depends
on \code{industrialio}, which has to be loaded first, as \file{run.sh} does with \code{modprobe}.

\xpart{d} A device that exchanges messages rather than producing values: an FPGA block that takes a
job description and returns a result of varying length, for instance. IIO's model is a set of
channels, each a value of fixed size, sampled together into scans of fixed size. Commands, replies
and records of varying length have no channel to be, so the driver would invent channels, spread
each message across them and leave userspace to put it back together, where a misc device that
reads and writes whole messages says what it does.

\xpart{e} That the claim is false. The GPIO character device needed a second version because the
first got things wrong, and the first is still built and supported on this kernel, marked
deprecated (\code{CONFIG_GPIO_CDEV_V1=y}; \file{drivers/gpio/Kconfig}:86--94), because programs use
it. Accurately: \textbf{a framework does not remove the ABI problem, it moves it to the subsystem's
maintainers}, who design, document and support the interface and its migrations once for every
driver in it, instead of every driver author doing it alone (\secref{c11:app:a8}).

% ------------------------------------------------------------------------------------------
\solution[fillaccent]{c11:ex:9}{A sample rate measured through an interface that knows nothing about it}{Cross-check}

\xpart{a} A \code{PERIOD_NS} of 1\,000\,000 is 1\,ms, so 1,000 samples a second, the first of them
one period after the device is enabled (\code{qa_dev_arm_timer}, \file{tools/qa-dev/qa-dev.c}:154).
In counter mode sample $n$ has the value $n$ and the count starts at zero (\file{qa-dev.c}:139--145),
so the thousandth sample has the value 999:
\begin{itemize}
  \item one second after the driver loads and enables the device, \code{in_voltage0_raw} reads
    \textbf{999}, or 998 if you read just before the thousandth sample is due;
  \item two reads a second apart differ by \textbf{1000};
  \item at \code{period_ns=500000}, 2,000 samples a second: \textbf{1999} and \textbf{2000}.
\end{itemize}
The first prediction assumes the count is at zero when the driver loads: after a boot in which the
device has not run, or with a driver that writes \code{CTRL_RESET}, which clears it
(\file{qa-dev.c}:268--272). Disabling the device stops its timer but keeps the count, so on a second
load in the same boot the value carries on from where it stopped, and only the difference can be
predicted.

\xpart{c} Reconcile:
\begin{itemize}
  \item About 880: \secref{c11:app:a4} measured 881, from 195 to 1076 in a second, 11.9 per cent
    short of 1,000. You have seen it in Chapter~\ref{c8:ch}'s Cross-check,
    Exercise~\ref{c8:ex:9}, as an interrupt count short of the period's prediction by a similar
    margin: 1,731 in about 2.1\,s at 1\,ms, which is about 824 a second. The effect is the device's
    timer. \code{qa-dev} arms the next sample one period after the moment its timer callback runs,
    not on an absolute schedule (\file{qa-dev.c}:154), so every interval is the period plus however
    late QEMU ran the callback. Exercise~\ref{c8:ex:9} also showed that this relative re-arm is
    part of the explanation and not all of it, because the lateness it implies is not constant.
  \item The same shortfall. The value is generated inside the device: the count goes up for every
    sample the device takes, even one it then drops because its FIFO is full (\file{qa-dev.c}:145
    and 173). So the difference between two values counts the samples the device produced, and
    nothing the driver or IIO does can add one or lose one; IIO decides only when you look. To tell,
    count something else over the same interval and compare: the change in \code{qa-dev}'s line of
    \file{/proc/interrupts}, which at 1\,ms is one interrupt per sample (Exercise~\ref{c8:ex:9}),
    should equal the change in the value. The one thing the reading path does change is the
    interval. From a shell, the time between two reads is the \code{sleep} plus the cost of
    starting the programs in between, which \secref{c11:app:a8} puts at tens of milliseconds a
    command. That makes a shell's figure a few per cent too high, never too low, so it cannot be the
    shortfall.
  \item Counter mode made it possible: each value is the device's own sample number, so the value
    is a running total, and you supplied the time yourself by reading a second apart. In
    \code{QA_DEV_MODE_NOISE} the values are a pseudo-random 12-bit sequence
    (\file{qa-dev.c}:133--136), and a difference means nothing. You would have had to count samples
    some other way: the interrupt count in \file{/proc/interrupts}, a count the driver publishes as
    Chapter~\ref{c10:ch}'s \code{samples} attribute did, or IIO's buffered mode with timestamps. Or
    put the device in sawtooth mode, where the value is $n$ modulo 4096 and a difference taken
    modulo 4096 still counts samples, as long as fewer than 4,096 arrive between two reads; 881 do.
\end{itemize}

\xpart{d} Measured on the course's target (\secref{c11:app:a8}): 571 microseconds a read from a C
program, 1,748 reads a second, against 29,700 microseconds from a shell loop of \code{cat}, 34 a
second. The C program is faster by $29\,700 / 571 \approx 52$ times, the chapter's ``fifty''. The
shell pays for a fork and an exec on every value, and even the C program pays for at least one
system call and for turning the number into text, for one value. The sysfs path is a way to
inspect a device, not a way to move samples.

\xpart{e} To see every sample, the reader must read at least as often as the device produces one,
so the limit is the rate of reads. From a shell that is about 34 a second: above about 30\,Hz a shell
loop misses samples (\secref{c11:app:a8}). From C it is 1,748 a second at best, before the program
does anything with the data; \code{qa-dev}'s 881 a second already takes half of that, so around
1\,kHz a loop is marginal, and above about 1.7\,kHz it cannot keep up at all. What IIO offers
instead is its buffered mode, triggered or not: the driver pushes each sample into a buffer in the
kernel as it arrives, and userspace reads them in batches from \file{/dev/iio:deviceN}, blocking
until they are there, so one system call carries many samples and none is lost while the buffer has
room. It exists because a value-at-a-time interface cannot keep up with a device that samples.

\xpart{f} ``Programmed for a 1\,ms period, 1,000 samples a second, \code{qa-dev} produced about 881
samples in a second, which is the difference between two counter values read through
\code{in_voltage0_raw} a second apart. The shortfall is in the device model's timing, which arms each
sample one period after its timer actually fired and is serviced late by QEMU, and not in the driver
or in IIO; in direct mode the driver delivers no stream at all, only the latest value when asked.''
Measured: the two values, their difference, and the interrupt count over the same interval if you
took it. Inferred: that the shortfall belongs to the device's timer rather than to the driver, which
rests on reading \file{qa-dev.c} and on Chapter~\ref{c8:ch}'s measurements; and that the interval
between the reads was one second, when from a shell it was a little longer.
